\documentclass[11pt]{article}
\usepackage{varnernotes}

\newcommand{\R}{\mathbb{R}}

\begin{document}

\lecturenotes{9a}{European Option Pricing with Black--Scholes--Merton}{Fall 2026}{October 20, 2026}{Jeffrey D. Varner}

\learningobjectives
  {Explain risk-neutral valuation}
  {Distinguish a no-arbitrage pricing measure from a forecast of the underlying's real-world return.}
  {Price European options}
  {Evaluate the Black--Scholes--Merton call and put formulas using consistent rate, volatility, strike, spot, dividend, and maturity conventions.}
  {Validate model premiums}
  {Use put--call parity, nonnegativity, strike monotonicity, and Monte Carlo standard errors as implementation checks.}

\section{European Exercise Reduces Valuation to a Terminal Payoff}

A European option can be exercised only at expiration $T$. Let $S_T\geq0$ be the underlying price in USD/share at expiration and let $K>0$ be the strike in USD/share. The per-share payoffs are
\begin{equation}
  \label{eq:payoffs}
  h_c(S_T;K):=(S_T-K)^+,
  \qquad
  h_p(S_T;K):=(K-S_T)^+,
  \qquad x^+:=\max\{x,0\}.
\end{equation}
The time-$0$ call and put premiums are denoted $C_0$ and $P_0$, respectively, and are measured in USD/share. The contract deliverable converts a per-share premium to account dollars; it does not change the per-share pricing formulas.

The exercise convention is load-bearing. The formulas in this lecture value the terminal payoffs in \eqnref{payoffs}; they do not optimize an early-exercise decision. L9b adds that decision and uses the European price as the comparison benchmark.

\section{Risk-Neutral Valuation Is Not a Return Forecast}

Let $S_t>0$ be the underlying price at time $t\in[0,T]$, let $g_y\in\R$ be the continuously compounded risk-free rate in the notation of L2a, let $\delta\in\R$ be a continuous proportional dividend yield, and let $\sigma>0$ be annualized volatility. Under the Black--Scholes--Merton (BSM) risk-neutral model,
\begin{equation}
  \label{eq:risk-neutral-gbm}
  \frac{dS_t}{S_t}=(g_y-\delta)\,dt+\sigma\,dW_t^{\mathbb Q},
\end{equation}
where $W^{\mathbb Q}$ is a Wiener process under the pricing measure $\mathbb Q$. The drift $g_y-\delta$ enforces the model's no-arbitrage pricing relation; it is not an estimate of the asset's historical or expected real-world return.

For a European payoff $h(S_T)$ with suitable integrability, the time-$0$ model value is
\begin{equation}
  \label{eq:risk-neutral-price}
  V_0=e^{-g_yT}\mathbb E^{\mathbb Q}[h(S_T)].
\end{equation}

\begin{remark}{What the model assumes}{bsm-assumptions}
The classical BSM derivation assumes frictionless continuous trading, no arbitrage, constant $g_y$, $\delta$, and $\sigma$, continuous lognormal price paths, and the ability to trade the underlying and financing asset as required by the hedge. Real markets have jumps, discrete trading, volatility smiles, funding spreads, transaction costs, and position constraints. A formula can be evaluated exactly while its model remains approximate.
\end{remark}

\section{Black--Scholes--Merton Closed Forms}

Let $S_0>0$, $K>0$, and $T>0$. Let $\Phi$ denote the standard-normal cumulative distribution function and define
\begin{equation}
  \label{eq:d1-d2}
  d_1=\frac{\ln(S_0/K)+(g_y-\delta+\tfrac12\sigma^2)T}{\sigma\sqrt T},
  \qquad
  d_2=d_1-\sigma\sqrt T.
\end{equation}

\begin{theorem}{BSM European call and put prices}{bsm}
Under \eqnref{risk-neutral-gbm} and the assumptions in \remref{bsm-assumptions}, the time-$0$ European premiums are
\begin{equation}
  \label{eq:bsm-prices}
  C_0=S_0e^{-\delta T}\Phi(d_1)-Ke^{-g_yT}\Phi(d_2),
  \qquad
  P_0=Ke^{-g_yT}\Phi(-d_2)-S_0e^{-\delta T}\Phi(-d_1).
\end{equation}
\end{theorem}

Every convention must agree across the inputs: $T$ is in years, $g_y$ is continuously compounded, $\sigma$ is annualized, and $S_0$ and $K$ use the same currency-per-share units. When the underlying pays no dividends, set $\delta=0$.

\subsection{Put--Call Parity}

The two European prices satisfy the replication identity
\begin{equation}
  \label{eq:put-call-parity}
  C_0-P_0=S_0e^{-\delta T}-Ke^{-g_yT}.
\end{equation}
Equation~\eqref{eq:put-call-parity} is an inexpensive implementation check. Compute the call and put directly and then test parity. Using a rounded call premium to obtain a tiny put by subtraction can suffer catastrophic cancellation and can even produce a negative numerical result.

\section{Monte Carlo Provides a Second Numerical Route}

Monte Carlo evaluates \eqnref{risk-neutral-price} by drawing $M$ independent risk-neutral endpoints, computing discounted payoffs $Y_k:=e^{-g_yT}h(S_T^{(k)})$, and using
\begin{equation}
  \label{eq:mc-estimator}
  \widehat V_M=\frac1M\sum_{k=1}^{M}Y_k,
  \qquad
  \widehat{\operatorname{SE}}(\widehat V_M)=\frac{s_Y}{\sqrt M},
\end{equation}
where $s_Y$ is the sample standard deviation. The closed form and simulation are different numerical routes to the same model value. Their difference should be interpreted in standard-error units; a nominal 95\% interval is not guaranteed to cover the closed form on every random run.

\notebooklink
  {L9a: European BSM and Risk-Neutral Monte Carlo}
  {https://github.com/varnerlab/CHEME-5660-CourseRepository-Fall-2026/blob/main/lectures/week-9/L9a/CHEME-5660-L9a-Example-BSM-Premium-Fall-2026.ipynb}
  {Price European calls and puts with both methods, quantify simulation error, verify put--call parity, and study premiums across strikes.}

\section{Checks Before Trusting a Premium}

Require nonnegative premiums. Under fixed remaining inputs, call premiums are nonincreasing in strike and put premiums are nondecreasing in strike. Verify put--call parity within the rounding tolerance of the implementation. For simulation, report the seed policy, path count, standard error, and standardized discrepancy from the closed form rather than only a point estimate.

These checks establish internal consistency under the selected model; they do not prove that the assumptions describe a particular market well. A disagreement becomes diagnostic evidence only after rate, dividend, volatility, time, exercise style, units, and timestamps have been reconciled.

\keytakeaways
  {In this lecture, we priced European calls and puts by risk-neutral valuation and the BSM closed forms, then checked the results with parity and simulation.}
  {Risk neutral is a pricing measure}
  {The drift under $\mathbb Q$ enforces no-arbitrage valuation; it is not a forecast of the asset's real-world return.}
  {BSM prices European exercise}
  {The call and put formulas value terminal payoffs under explicit assumptions and consistent input conventions.}
  {A model premium needs checks}
  {Direct formulas, put--call parity, strike monotonicity, and Monte Carlo standard errors reveal many implementation and convention errors.}
  {Next, add the American exercise right and measure when that flexibility makes the American premium exceed the European benchmark.}

\begin{readinglist}
  \item Fischer Black and Myron Scholes, \href{https://doi.org/10.1086/260062}{``The Pricing of Options and Corporate Liabilities,''} \emph{Journal of Political Economy} 81(3), 637--654 (1973).
  \item Robert C. Merton, \href{https://doi.org/10.2307/1831029}{``Theory of Rational Option Pricing,''} \emph{Bell Journal of Economics and Management Science} 4(1), 141--183 (1973).
\end{readinglist}

\vfill
{\sffamily\footnotesize\color{vnmuted}\textbf{Educational use and risk notice.} These notes are for instruction only and do not constitute investment advice. Option values are model-dependent and trading involves substantial risk.}

\end{document}
